\section{置換後の計算}\label{節:置換後の整理}

置換を選ぶとき,「すべてがすぐ基本公式になるか」だけを条件にすると,有用な置換を見落とす.根号をなくす,指数の中身を一つにする,分母の二次式を一次式にするなど,まず何を簡単にしたいかを決める.その後で,前章の式の整理を使えばよい.

\begin{example}{置換後の整理}
$\int\frac{x^3}{1+x^2}\dd x$ を求めよ.
\end{example}

分母の $1+x^2$ を $t$ と置くと,\,$\dd t=2x\dd x$ である.分子には $x^3$ があるが,\,$x^3\dd x=x^2(x\dd x)$ と分ければ,\,$x\dd x=\frac12\dd t$ を使える.残る $x^2$ も $t-1$ と書けるので,
\[
\begin{aligned}
 \int\frac{x^3}{1+x^2}\dd x
 &=\frac12\int\frac{t-1}{t}\dd t\\
 &=\frac12\int\left(1-\frac1t\right)\dd t
 =\frac12t-\frac12\log t+C.
\end{aligned}
\]
$t=1+x^2>0$ を戻し,定数 $\frac12$ を任意定数へまとめると
\[
 \int\frac{x^3}{1+x^2}\dd x
 =\frac{x^2}{2}-\frac12\log(1+x^2)+C.
\]
$t=x^2+1$ は全実数上で一対一ではないが,残った $x^2$ を $t-1$ と書くのに逆関数 $x=h(t)$ は使っていない.\,$x$ の符号を選ぶ必要もない.

別の道もある.最初に割り算をすれば
\[
 \frac{x^3}{1+x^2}=x-\frac{x}{1+x^2}
\]
なので,第一項をそのまま積分し,第二項だけに置換を使って同じ答えを得る.どちらの場合も,式のどこを簡単にするかが分かっていれば,計算の途中で方針を見失わない.

これまでの置換では,複雑なまとまりとその微分が結び付いていた.しかし $x\e^x$ のように,二つの因子を片方の合成関数の微分としてまとめにくい積もある.この場合には,積の微分を逆に読む方法が役立つ.次章では,積分する因子と微分する因子を分けて考える.


